\subsection{The naturality class of embedded smooth partitions}
\label{subsec:embedded-smooth-naturality}

\begin{proposition}[Smooth-locus preservation and the cusp]
\label{prop:embedded-smooth-partition-naturality}
Let $X\subset\mathbb R^N$ and $X'\subset\mathbb R^{N'}$ be
semialgebraic sets, with finite marked families $\Phi$ and $\Phi'$.
A semialgebraic homeomorphism $h:X\to X'$ that carries the marks to
the corresponding marks and satisfies
\[
 h(\Sigma(Y))=\Sigma(h(Y))
 \quad\text{for every semialgebraic }Y\subset X
 \text{ closed in }X
\]
carries each of the coarsest partitions in
Theorems~\ref{thm:coarsest-embedded-smooth-partition}
and~\ref{thm:coarsest-smooth-partition-variants} to its counterpart
on $X'$.  In particular this holds for the restrictions of ambient
Nash diffeomorphisms.  It fails for arbitrary semialgebraic
homeomorphisms, already for the interval and the compact cusp.
\end{proposition}

\begin{proof}
A homeomorphism transports relative closures, local closedness,
connectedness, and the frontier condition.  A semialgebraic
homeomorphism also preserves local dimension, by applying
semialgebraic invariance of dimension to corresponding neighborhoods
\cite[Proposition~1.8, p.~5]{Coste2007}.  The displayed assumption
therefore transports the $F$-regularity condition for every closed
saturated union.  Its inverse has the same property, because every
closed semialgebraic subset of $X'$ is the image of one in $X$.
Thus transport gives mutually inverse, order-preserving bijections
between the relevant classes of adapted partitions.  It carries
their unique coarsest elements to each other.  Preservation of the
smooth-locus operator by Nash diffeomorphisms is precisely
\cite[Lemma~2.3, p.~636]{Fernando2018}.

For failure of general homeomorphism invariance, let
\[
 I=[-1,1],\qquad
 C=\{(t^2,t^3):-1\le t\le1\},\qquad
 \gamma(t)=(t^2,t^3).
\]
The map $\gamma:I\to C$ is a semialgebraic homeomorphism.
It is bijective, and its inverse is $y/x$ away from the origin,
extended by zero at the origin.  Continuity there follows from
$|\gamma^{-1}(x,y)|=\sqrt{x}$ on $C$.
The origin is an embedded smooth point of $I$, but not of $C$.
Indeed $C$ has local dimension one at the origin, and its tangent
cone there is the nonnegative horizontal half-line.  For example
\[
 \frac{\gamma(t)}{\|\gamma(t)\|}
   =\frac{(1,t)}{\sqrt{1+t^2}}
   \longrightarrow(1,0)
 \quad (t\longrightarrow0,\ t\ne0).
\]
Every nonzero tangent-cone direction is therefore positive horizontal,
and every positive horizontal vector occurs by rescaling.
An embedded $C^1$ one-manifold without boundary has a full tangent
line, so this half-line excludes such a manifold germ.

The nonsmooth loci of the two compact sets are exactly
\[
 \Sigma(I)=\{-1,1\},\qquad
 \Sigma(C)=\{\gamma(-1),(0,0),\gamma(1)\}.
\]
The endpoints have one-sided manifold germs and hence are nonsmooth
under the convention without boundary.  Away from those endpoints
and the cusp origin, $\gamma$ is a regular Nash parametrization of
an embedded open arc.

In either case the partition into the displayed finite nonsmooth
set and its complement is pure-dimensional, cellular, and
$F$-regular.  Its only closed saturated unions are the empty set,
the finite zero-dimensional stratum, and the whole space; their
nonsmooth loci are respectively empty, empty, and that finite
stratum.  Every admissible partition must saturate the nonsmooth
locus of the whole space, so it refines this two-piece partition.
These are therefore the coarsest pure-dimensional partitions.
The open interval is one stratum, while the cusp origin belongs
to a zero-dimensional stratum containing both endpoints as well.
Since the finite strata have two and three
points respectively, no homeomorphism can identify the two
partitions.  Separating the finite points and the connected open arcs
gives admissible partitions with respectively three and five strata:
every closed saturated union is a union of closed subarcs and isolated
vertices, and its nonsmooth locus is a union of vertices.  A connected
partition refining either two-piece partition also refines its
component partition.  Hence these are the two coarsest connected
versions.
\end{proof}

